\scanpage{90}
\mspage{\handwritten{65}}

considerably wider sense than I expected, when writing this letter to Larry Breen. It now appears that the whole theory of stacks of groupoids will depend on the choice of a ``coherator'' $\underline{\mathrm{C}}$ , as seen in section 13.

(\textsuperscript{26}) Tim Porter pointed out to me that ``Dold Puppe'' is an accurate name for this basic theorem, which should be called Dold-Kan theorem.

(\textsuperscript{27}) As was pointe\deleted{\unsure{e}}d out to me by Ronnie Brown, this structure was already well-known to J.H.C. Whitehead, under the name of ``crossed modules'', and extensive use and extensive generalizations of this notion (in quite different directions from those I was having in mind, in terms of Gr-stacks over an arbitrary topos) have been made by him and others. With respect to the question on next page, of generalizing this notion of ``non commuative chain complex'' from length one to length two, Ronnie says there is a work in preparation by D. Conduché `` Modules croisés généralisés de longueur 2 '' .

(\textsuperscript{28}) Tim P\deleted{\unsure{e}}orter pointed out to me that work on étale tubular neighbourhoods was done by D.A. Cox ; ``Algebraic tubular neighbourhoods I,II', Math. Scand. 42 (1978) 211-228, 229-242. I've not seen yet this work, and can't say therefore whether it meets the rather precise expectations I have for a theory of tubular neighbourhoods, for the needs of a dévissage theory of stratified schemes (or, more generally, stratified topoi).

(\textsuperscript{29}) This is the typical game \deleted{to be} embodied in the ``derivator'' associa\deleted{\unsure{te}}ted to the theory $(\mathrm{Hot})$ of usual homotopy types (compare section 69).

(\textsuperscript{30}) Some more details on this program are outlined in Esquisse d'un \deleted{\unsure{P}}programme (section 5), in Réflexions Mathématiques 1 .

(\textsuperscript{31}) I was informed by knowledgeable people soon later that the answer is well known to be negative, by working with ``rational homotopy types'' (the cohomology of which is made up with vector spaces over $\underline{\mathrm{Q}}$ . It is well known indeed that a 1-connected rational homotopy type is not known from its rational cohomology ring alone, which is supposed (I confess I did'nt check) to contain already all the information I was contemplating.

(\textsuperscript{32}) This \deleted{is the \unsure{``now standard''|now standard} version of the similar problem}\inserted{``problem'' is met with by the notion of a ``derivator''} , which ``was in the air'' already by the late sixties, but was never developped - (instead even derived categories became tabu in the seventies ...). Compare with note (\textsuperscript{24}).
% archive sheet number: 89
